\documentclass[11pt,a4paper]{article}
\usepackage[T1]{fontenc}
\usepackage{lmodern,amsmath,amssymb,amsthm,mathtools,booktabs,microtype}
\usepackage[margin=26mm]{geometry}
\usepackage[colorlinks=true,linkcolor=blue,citecolor=blue,urlcolor=blue]{hyperref}
\setlength{\emergencystretch}{2em}
\newtheorem{theorem}{Theorem}[section]
\newtheorem{proposition}[theorem]{Proposition}
\newtheorem{lemma}[theorem]{Lemma}
\newtheorem{corollary}[theorem]{Corollary}
\theoremstyle{definition}\newtheorem{definition}[theorem]{Definition}
\theoremstyle{remark}\newtheorem{remark}[theorem]{Remark}
\newcommand{\Q}{\mathbb Q}\newcommand{\Z}{\mathbb Z}
\newcommand{\R}{\mathbb R}\newcommand{\A}{\mathbb A}
\newcommand{\PP}{\mathbb P}\newcommand{\sqcl}[1]{[#1]}
\DeclareMathOperator{\Tr}{Tr}\DeclareMathOperator{\Nm}{N}
\DeclareMathOperator{\N}{N}\DeclareMathOperator{\rad}{rad}
\DeclareMathOperator{\sfpart}{sf}
\title{Rational Local Approximations to Integral Pell Labels}
\author{}\date{16 November 2026}
\begin{document}\maketitle
\begin{abstract}
For every positive integer $r$, the equation
$t^2-5tw+w^2=1$, with $t=3n+1$, has a rational solution with
$0<n<1$ and $v_{19}(n-1)\ge r$. The exact integral target
$n=1$ has no rational solution. A rational parametrization and
a simple Hensel lift prove the assertion for every precision.
Thus approximation at this prime cannot replace the exact
integrality premise in the Pell-label growth argument.
\end{abstract}

\section{Norm equations and exact labels}
Consider
\begin{equation}\label{eq:norm}
                    t^2-atw+w^2=1,\qquad t=(a-2)n+1.
\end{equation}
When $a,t$ are integers, a rational solution $w$ is integral by
monicity. Under the further hypotheses stated in~\cite{IntegralLabels},
the congruence $t\equiv1\pmod{|a-2|}$ forces a sharp exponential
lower bound on the nonzero integer label $n$. The following explicit
example keeps $a=5$ and the norm equation exact but replaces the
integral label by a rational number. We normalize $v_{19}$ by
$v_{19}(19)=1$.

\begin{proposition}\label{prop:local}
For every integer $r\ge1$ there are rational numbers $n_r,w_r$ such
that, with $t_r=3n_r+1$,
\[
 t_r^2-5t_rw_r+w_r^2=1,\qquad
 v_{19}(n_r-1)\ge r,\qquad 0<n_r<1.
\]
In contrast the exact integral labels $a=5,n=1$ admit no rational
solution of the norm equation.
\end{proposition}
\begin{proof}
The integer $85$ has the simple square root $3$ modulo $19$.
Lift it to an integer $q_r$ satisfying
$q_r^2\equiv85\pmod{19^r}$ and $q_r\equiv3\pmod{19}$.
Choose $s_r\ge19$ with
\[
                       3s_r\equiv10+q_r\pmod{19^r}.
\]
Put
\begin{equation}\label{eq:param}
 D_r=s_r^2-5s_r+1,\qquad
 t_r=\frac{s_r^2-1}{D_r},\qquad
 w_r=\frac{s_r(5s_r-2)}{D_r},\qquad
 n_r=\frac{t_r-1}{3}.
\end{equation}
Substitution proves the norm equation identically in $s_r$.
Modulo $19$, one has $s_r\equiv17$ and $D_r\equiv15$,
so every denominator in~\eqref{eq:param} is a $19$-adic unit.
Moreover
\[
 3D_r(t_r-4)=-9s_r^2+60s_r-15
                  =85-(3s_r-10)^2.
\]
It follows that $v_{19}(t_r-4)\ge r$, and hence
$v_{19}(n_r-1)\ge r$. Since $s_r\ge19$, one has $D_r>0$,
$t_r>1$ from $5s_r-2>0$, and $t_r<4$ from
$3s_r^2-20s_r+5>0$. Thus $0<n_r<1$.

For the integral target $a=5,n=1$, the value of $t$ is four and
the discriminant is $340=4\cdot85$, which is not a rational square.
The quadratic formula therefore rules out a rational value of $w$.
\end{proof}


The proposition concerns arbitrarily accurate approximation at the
single specified prime $19$. It shows that this local requirement,
even together with the exact norm equation, cannot replace exact
integral labels. It does not claim such a counterexample for every
finite set of places or for additional arithmetic restrictions.

\section{Exact verification}
The standard-library program
\nolinkurl{computation/integer_labels/pell_compression_v40/check.py}
checks the parametrization identity and explicit local approximants
with exact rational arithmetic. It also checks the identities and
edge cases of the integral-label theorem~\cite{IntegralLabels}.
The proof above establishes the assertion for every $r$; it is not
inferred from a finite collection of approximants.

\begin{thebibliography}{99}
\bibitem{IntegralLabels}
\emph{Integral Lucas labels from a finite quadratic choice},
project manuscript, 15 November 2026,
\nolinkurl{papers/integer_labels/pell_labels_20261115_v40.tex}.
\end{thebibliography}
\end{document}
